\def\EMtitlesize{24}
\begin{frame}{Part I}
\begin{tikzpicture}[remember picture,overlay]
\node[fill=SLbar,draw=SLbarEdge,line width=0.8bp,rounded corners=4bp,minimum width=520bp,minimum height=170bp] (b) at ([yshift=-20bp]current page.center) {};
\node[text=SLthm,font=\bfseries\fontsize{30}{36}\selectfont] at ([yshift=36bp]b.center) {\EMdir{Part I}};
\node[text=SLdef,font=\fontsize{24}{30}\selectfont] at ([yshift=-30bp]b.center) {\EMdir{Fourier Analysis}};
\end{tikzpicture}
\end{frame}

\section{Fourier Series}
\def\EMfootchapter{Chapter 1 \ \textbullet\ Fourier Series}
\def\EMtitlesize{24}
\begin{frame}{Chapter 1: Fourier Series}
{\fontsize{20}{26}\selectfont\bfseries\color{SLdef}\EMalign In this chapter\par}
\vspace{10bp}
\EMitem{\textbullet}{\dimexpr\linewidth-26bp\relax}{Periodic Functions}
\EMitem{\textbullet}{\dimexpr\linewidth-26bp\relax}{Review of Trigonometry}
\EMitem{\textbullet}{\dimexpr\linewidth-26bp\relax}{Review: Integrals and Orthogonality}
\EMitem{\textbullet}{\dimexpr\linewidth-26bp\relax}{The Fourier Series}
\EMitem{\textbullet}{\dimexpr\linewidth-26bp\relax}{Review: The Definite Integral of a Product of Two Functions}
\EMitem{\textbullet}{\dimexpr\linewidth-26bp\relax}{Examples of Fourier Series}
\EMitem{\textbullet}{\dimexpr\linewidth-26bp\relax}{Fourier Series of Functions with an Arbitrary Period}
\EMitem{\textbullet}{\dimexpr\linewidth-26bp\relax}{Dirichlet's Theorem}
\EMitem{\textbullet}{\dimexpr\linewidth-26bp\relax}{Solving Some Problems}
\end{frame}

\def\EMtitlesize{24}
\begin{frame}{Periodic Functions}
\begin{minipage}[t]{500bp}\vspace{0pt}
\begin{uncoverenv}<2->
\EMeq{f(x+p) =f(x)\quad \Longrightarrow\quad f(x+np) =f(x)}
\end{uncoverenv}
\EMgap
\EMgap
\begin{uncoverenv}<4->
\EMitem{\textbullet}{\dimexpr\linewidth-26bp\relax}{The smallest positive value of \EMim{p} for which the function is periodic is called the fundamental period of the function.}
\end{uncoverenv}
\EMgap
\begin{uncoverenv}<5->
\EMitem{\textbullet}{\dimexpr\linewidth-26bp\relax}{If a function is periodic with period \EMim{p}, it is also periodic with period \EMim{np}, where \EMim{n} is a natural number.}
\end{uncoverenv}
\end{minipage}
\EMside{36.1bp}{\uncover<3->{\EMfigure{1.140}{m01-periodic}{A periodic function with period \EMim{p}}}}
\end{frame}

\def\EMtitlesize{24}
\begin{frame}{The Periodic Functions Sine and Cosine}
\begin{minipage}[t]{500bp}\vspace{0pt}
\begin{EMbox}{ex}{1}{\textbf{Example 1.1}}
\EMp{The function \EMim{\sin x} is periodic with period \EMim{2\pi}, and \EMim{\sin2x} is periodic with fundamental period \EMim{\pi}, although \EMim{2\pi} is also a period of it. Likewise \EMim{\sin3x} is periodic with period \EMim{\frac{2\pi}{3}}, although \EMim{2\pi} is also a period of it.}
\end{EMbox}
\end{minipage}
\EMside{145.1bp}{\uncover<2->{\EMfigure{1.023}{m01-sincos}{The functions \EMim{\sin nx} and \EMim{\cos nx} for \EMim{n=1,2,3} on \EMim{[0,2\pi]}}}}
\end{frame}

\def\EMtitlesize{24}
\begin{frame}{Review of Trigonometry}
\begin{uncoverenv}<2->
\EMeq{\begin{gathered}
\sin^2\alpha+\cos^2\alpha=1\\
\sin(\alpha+\beta)=\sin\alpha\cos\beta+\cos\alpha\sin\beta\\
\cos(\alpha+\beta)=\cos\alpha\cos\beta-\sin\alpha\sin\beta\\
\cos2\alpha=2\cos^2\alpha-1=1-2\sin^2\alpha=\cos^2\alpha-\sin^2\alpha\\
\sin2\alpha=2\sin\alpha\cos\alpha\\
\sin^2\alpha=\frac{1-\cos2\alpha}{2}\qquad\cos^2\alpha=\frac{1+\cos2\alpha}{2}\\
\sin\alpha\sin\beta=\tfrac12\big[\cos(\alpha-\beta)-\cos(\alpha+\beta)\big]\\
\cos\alpha\cos\beta=\tfrac12\big[\cos(\alpha-\beta)+\cos(\alpha+\beta)\big]\\
\sin\alpha\cos\beta=\tfrac12\big[\sin(\alpha-\beta)+\sin(\alpha+\beta)\big]
\end{gathered}}
\end{uncoverenv}
\end{frame}

\def\EMtitlesize{24}
\begin{frame}{Review: Integrals and Orthogonality}
\begin{uncoverenv}<2->
\EMeq{\int_{-\pi}^{\pi}\sin nt\,dt =0,\qquad n\in\mathbb{Z}}
\end{uncoverenv}
\EMgap
\begin{uncoverenv}<3->
\EMeq{\int_{-\pi}^{\pi}\cos nt\,dt =0,\qquad n\in\mathbb{Z}\setminus\{0\}}
\end{uncoverenv}
\EMgap
\begin{EMbox}{prf}{3}{\textbf{Proof}}
\begin{uncoverenv}<4->
\EMeq{\int_{-\pi}^{\pi}\sin nt\,dt =-\frac1n\cos nt\Big|_{-\pi}^{\pi} =0,\qquad \int_{-\pi}^{\pi}\cos nt\,dt =\frac1n\sin nt\Big|_{-\pi}^{\pi} =0\qquad (n\neq0)}
\end{uncoverenv}
\end{EMbox}
\end{frame}

\def\EMtitlesize{24}
\begin{frame}{Orthogonality}
\begin{uncoverenv}<2->
\EMeq{\int_{-\pi}^{\pi}\sin nt\sin mt\,dt =0,\qquad n,m\in\mathbb{Z},\ n\neq m}
\end{uncoverenv}
\EMgap
\begin{uncoverenv}<3->
\EMeq{\int_{-\pi}^{\pi}\cos nt\cos mt\,dt =0,\qquad n,m\in\mathbb{Z},\ n\neq m}
\end{uncoverenv}
\EMgap
\begin{uncoverenv}<4->
\EMeq{\int_{-\pi}^{\pi}\sin nt\cos mt\,dt =0,\qquad n,m\in\mathbb{Z}\ \ (n\neq m\ \text{or}\ n =m)}
\end{uncoverenv}
\EMgap
\begin{EMbox}{prf}{4}{\textbf{Proof}}
\begin{uncoverenv}<5->
\EMeq{\int_{-\pi}^{\pi}\sin nx\sin mx\,dx =\frac12\int_{-\pi}^{\pi}\cos(n-m)x\,dx-\frac12\int_{-\pi}^{\pi}\cos(n+m)x\,dx}
\end{uncoverenv}
\EMgap
\begin{uncoverenv}<6->
\EMeq{\int_{-\pi}^{\pi}\cos nx\cos mx\,dx =\frac12\int_{-\pi}^{\pi}\cos(n+m)x\,dx+\frac12\int_{-\pi}^{\pi}\cos(n-m)x\,dx}
\end{uncoverenv}
\EMgap
\begin{uncoverenv}<7->
\EMeq{\int_{-\pi}^{\pi}\sin nx\cos mx\,dx =\frac12\int_{-\pi}^{\pi}\sin(n+m)x\,dx+\frac12\int_{-\pi}^{\pi}\sin(n-m)x\,dx}
\end{uncoverenv}
\end{EMbox}
\end{frame}

\def\EMtitlesize{24}
\begin{frame}{Orthogonality (continued)}
\begin{uncoverenv}<2->
\EMeq{\int_{-\pi}^{\pi}\sin^2nt\,dt =\pi,\qquad \int_{-\pi}^{\pi}\cos^2nt\,dt =\pi,\qquad n\in\mathbb{Z}\setminus\{0\}}
\end{uncoverenv}
\EMgap
\begin{EMbox}{prf}{2}{\textbf{Proof}}
\begin{uncoverenv}<3->
\EMeq{\int_{-\pi}^{\pi}\sin^2nt\,dt =\int_{-\pi}^{\pi}\frac{1-\cos2nt}{2}\,dt =\pi,\qquad \int_{-\pi}^{\pi}\cos^2nt\,dt =\int_{-\pi}^{\pi}\frac{1+\cos2nt}{2}\,dt =\pi}
\end{uncoverenv}
\end{EMbox}
\end{frame}

\def\EMtitlesize{24}
\begin{frame}{The Fourier Series}
\begin{EMbox}{rem}{1}{\textbf{Remark}\ \ |\ \ History}
\begin{uncoverenv}<2->
\EMp{Joseph Fourier (1768--1830) put forward the idea of expressing periodic functions as a series of sinusoidal functions. When this revolutionary idea was presented, some mathematicians, Lagrange among them, objected to its generality, until Dirichlet stated the conditions for the convergence of the Fourier series.}
\end{uncoverenv}
\end{EMbox}
\begin{EMbox}{thm}{2}{\textbf{Theorem 1.2}\ \ |\ \ The Fourier series}
\begin{uncoverenv}<3->
\EMp{If the function \EMim{f(x)} is periodic with period \EMim{2\pi}, it can be expressed as a trigonometric series as follows:}
\vspace{3bp}
\EMeq{f(x) =a_0+\sum_{n=1}^{\infty}\big(a_n\cos nx+b_n\sin nx\big)}
\end{uncoverenv}
\EMgap
\begin{uncoverenv}<4->
\EMp{where}
\vspace{3bp}
\EMeq{a_0 =\frac1{2\pi}\int_{-\pi}^{\pi}f(x)\,dx,\qquad a_n =\frac1\pi\int_{-\pi}^{\pi}f(x)\cos nx\,dx,\qquad b_n =\frac1\pi\int_{-\pi}^{\pi}f(x)\sin nx\,dx}
\end{uncoverenv}
\end{EMbox}
\end{frame}

\def\EMtitlesize{24}
\begin{frame}{The Fourier Series (continued)}
\begin{EMbox}{prf}{1}{\textbf{Proof}}
\begin{uncoverenv}<2->
\EMp{Suppose the periodic function has been written as the trigonometric series above; it suffices to find the coefficients \EMim{a_0}, \EMim{a_n} and \EMim{b_n}.}
\end{uncoverenv}
\EMgap
\begin{uncoverenv}<3->
\EMp{\textbf{Computing \EMim{a_0}.} Integrating both sides over an interval of length \EMim{2\pi} gives}
\vspace{3bp}
\EMeq{\int_{-\pi}^{\pi}f(x)\,dx =\int_{-\pi}^{\pi}\left[a_0+\sum_{n=1}^{\infty}(a_n\cos nx+b_n\sin nx)\right]dx =a_0\int_{-\pi}^{\pi}dx+\sum_{n=1}^{\infty}\left(a_n\underbrace{\int_{-\pi}^{\pi}\cos nx\,dx}_{0}+b_n\underbrace{\int_{-\pi}^{\pi}\sin nx\,dx}_{0}\right) \\[4pt] =2\pi a_0}
\end{uncoverenv}
\EMgap
\begin{uncoverenv}<4->
\EMeq{\Longrightarrow\quad a_0 =\frac1{2\pi}\int_{-\pi}^{\pi}f(x)\,dx}
\end{uncoverenv}
\EMgap
\begin{uncoverenv}<5->
\EMp{\EMim{a_0} is the mean value or DC level of the function.}
\end{uncoverenv}
\end{EMbox}
\end{frame}

\def\EMtitlesize{24}
\begin{frame}{The Fourier Series (continued)}
\begin{EMbox}{prf}{1}{}
\begin{uncoverenv}<2->
\EMp{\textbf{Computing \EMim{a_n}.} Multiplying both sides by \EMim{\cos mx} and integrating over an interval of length \EMim{2\pi}:}
\vspace{3bp}
\EMeq{\int_{-\pi}^{\pi}f(x)\cos mx\,dx =\int_{-\pi}^{\pi}a_0\cos mx\,dx+\sum_{n=1}^{\infty}\int_{-\pi}^{\pi}a_n\cos nx\cos mx\,dx+\sum_{n=1}^{\infty}\int_{-\pi}^{\pi}b_n\sin nx\cos mx\,dx =\pi a_m}
\end{uncoverenv}
\EMgap
\begin{uncoverenv}<3->
\EMeq{\Longrightarrow\quad a_n =\frac1\pi\int_{-\pi}^{\pi}f(x)\cos nx\,dx}
\end{uncoverenv}
\EMgap
\begin{uncoverenv}<4->
\EMp{\textbf{Computing \EMim{b_n}.} Multiplying both sides by \EMim{\sin mx} and integrating over an interval of length \EMim{2\pi}:}
\vspace{3bp}
\EMeq{\int_{-\pi}^{\pi}f(x)\sin mx\,dx =\int_{-\pi}^{\pi}a_0\sin mx\,dx+\sum_{n=1}^{\infty}\int_{-\pi}^{\pi}a_n\cos nx\sin mx\,dx+\sum_{n=1}^{\infty}\int_{-\pi}^{\pi}b_n\sin nx\sin mx\,dx =\pi b_m}
\end{uncoverenv}
\EMgap
\begin{uncoverenv}<5->
\EMeq{\Longrightarrow\quad b_n =\frac1\pi\int_{-\pi}^{\pi}f(x)\sin nx\,dx}
\end{uncoverenv}
\end{EMbox}
\end{frame}

\def\EMtitlesize{24}
\begin{frame}{The Fourier Series (continued)}
\begin{EMbox}{imp}{1}{\textbf{Key result 1.3}\ \ |\ \ The Fourier series and Euler's formulas}
\begin{uncoverenv}<2->
\EMeq{f(x) =a_0+\sum_{n=1}^{\infty}\big(a_n\cos nx+b_n\sin nx\big),\qquad f(x+2\pi) =f(x)}
\end{uncoverenv}
\EMgap
\begin{uncoverenv}<3->
\EMeq{a_0 =\frac1{2\pi}\int_{2\pi}f(x)\,dx,\qquad a_n =\frac1\pi\int_{2\pi}f(x)\cos nx\,dx,\qquad b_n =\frac1\pi\int_{2\pi}f(x)\sin nx\,dx}
\end{uncoverenv}
\EMgap
\begin{uncoverenv}<4->
\EMp{(the integral is taken over one period of length \EMim{2\pi}.)}
\end{uncoverenv}
\end{EMbox}
\begin{EMbox}{rem}{4}{\textbf{Remark}}
\begin{uncoverenv}<5->
\EMp{At points of discontinuity of the function, the value of the Fourier series equals the average of the upper and lower limits at that point.}
\end{uncoverenv}
\end{EMbox}
\end{frame}

\def\EMtitlesize{24}
\begin{frame}{Review: The Definite Integral of a Product of Two Functions}
\begin{uncoverenv}<2->
\EMp{By integration by parts:}
\vspace{3bp}
\EMeq{\int_\alpha^\beta f(x)g(x)\,dx =\Big[f(x)G_1(x)-f^{(1)}(x)G_2(x)+f^{(2)}(x)G_3(x)-+\cdots\Big]_\alpha^\beta}
\end{uncoverenv}
\EMgap
\begin{uncoverenv}<3->
\EMp{where \EMim{f^{(k)}} is the \EMim{k}-th derivative of \EMim{f} and \EMim{G_k} is the \EMim{k}-th integral of \EMim{g}:}
\end{uncoverenv}
\EMgap
\begin{uncoverenv}<4->
{\EMbody\centering\scalebox{1.000}{\parbox{880bp}{\setlength{\columnwidth}{792pt}\centering \begin{tabular}{@{}ccc@{}}
\toprule\noalign{}
\rowcolor{EMindigoL}derivatives of \EMim{f} & & integrals of \EMim{g} \\\noalign{\vskip 4pt}
\midrule\noalign{}

\EMim{f(x)} & \EMim{\searrow} & \EMim{g(x)} \\\noalign{\vskip 4pt}
\EMim{f^{(1)}(x)} & \EMim{\searrow} & \EMim{G_1(x)} \\\noalign{\vskip 4pt}
\EMim{f^{(2)}(x)} & \EMim{\searrow} & \EMim{G_2(x)} \\\noalign{\vskip 4pt}
\EMim{f^{(3)}(x)} & \EMim{\searrow} & \EMim{G_3(x)} \\\noalign{\vskip 4pt}
\EMim{\vdots} & & \EMim{\vdots} \\\noalign{\vskip 4pt}
\bottomrule
\end{tabular}}}\par}
\end{uncoverenv}
\end{frame}

\def\EMtitlesize{24}
\begin{frame}{Examples of Fourier Series}
\begin{EMbox}{ex}{1}{\textbf{Example 1.4}}
\EMp{Find the Fourier series of the following function:}
\vspace{3bp}
\EMeq{f(x) =\begin{cases}-k&-\pi<x<0\\ \ \ k&0<x<\pi\end{cases}\qquad f(x+2\pi) =f(x)}
\EMgap
\begin{uncoverenv}<2->
\EMeq{a_0 =\frac1{2\pi}\int_{-\pi}^{\pi}f(x)\,dx =0}
\end{uncoverenv}
\EMgap
\begin{uncoverenv}<3->
\EMeq{a_n =\frac1\pi\int_{-\pi}^{\pi}f(x)\cos nx\,dx =\frac1\pi\left[\int_{-\pi}^{0}(-k)\cos nx\,dx+\int_0^\pi k\cos nx\,dx\right] =\frac1\pi\left[-k\,\frac{\sin nx}{n}\Big|_{-\pi}^{0}+k\,\frac{\sin nx}{n}\Big|_0^\pi\right] =0}
\end{uncoverenv}
\EMgap
\begin{uncoverenv}<4->
\EMeq{\begin{aligned}
b_n&=\frac1\pi\int_{-\pi}^{\pi}f(x)\sin nx\,dx=\frac1\pi\left[\int_{-\pi}^{0}(-k)\sin nx\,dx+\int_0^\pi k\sin nx\,dx\right]\\
&=\frac1\pi\left[k\,\frac{\cos nx}{n}\Big|_{-\pi}^{0}-k\,\frac{\cos nx}{n}\Big|_0^\pi\right]=\frac k{n\pi}\big[\cos0-\cos(-n\pi)-\cos n\pi+\cos0\big]=\frac{2k}{n\pi}(1-\cos n\pi)
\end{aligned}}
\end{uncoverenv}
\end{EMbox}
\end{frame}

\def\EMtitlesize{24}
\begin{frame}{Examples of Fourier Series (continued)}
\begin{minipage}[t]{500bp}\vspace{0pt}
\begin{EMbox}{ex}{1}{}
\begin{uncoverenv}<2->
\EMeq{\cos n\pi =\begin{cases}-1&n\ \text{odd}\\ \ \ 1&n\ \text{even}\end{cases}\quad \Longrightarrow\quad 1-\cos n\pi =\begin{cases}2&n\ \text{odd}\\ 0&n\ \text{even}\end{cases}}
\end{uncoverenv}
\EMgap
\begin{uncoverenv}<3->
\EMeq{b_1 =\frac{4k}{\pi},\quad b_2 =0,\quad b_3 =\frac{4k}{3\pi},\quad b_4 =0,\quad b_5 =\frac{4k}{5\pi},\ \dots}
\end{uncoverenv}
\EMgap
\begin{uncoverenv}<4->
\EMeq{f(x) =\frac{4k}\pi\left(\sin x+\frac13\sin3x+\frac15\sin5x+\cdots\right)}
\end{uncoverenv}
\end{EMbox}
\end{minipage}
\EMside{199.6bp}{\uncover<5->{\EMfigure{1.130}{m01-square-wave}{A square wave \EMim{f(x)} with period \EMim{2\pi} and values \EMim{\pm k}}}}
\end{frame}

\def\EMtitlesize{24}
\begin{frame}{Examples of Fourier Series (continued)}
\begin{minipage}[t]{500bp}\vspace{0pt}
\begin{uncoverenv}<2->
\EMp{The partial sums of the series:}
\end{uncoverenv}
\EMgap
\begin{uncoverenv}<3->
\EMeq{S_1 =\frac{4k}\pi\sin x,\qquad S_2 =\frac{4k}\pi\left(\sin x+\frac13\sin3x\right),\qquad S_3 \\[4pt] =\frac{4k}\pi\left(\sin x+\frac13\sin3x+\frac15\sin5x\right)}
\end{uncoverenv}
\EMgap
\EMgap
\begin{uncoverenv}<5->
\EMp{Setting \EMim{x=\frac\pi2}:}
\vspace{3bp}
\EMeq{f\!\left(\frac\pi2\right) =\frac{4k}\pi\left(1-\frac13+\frac15-\frac17+-\cdots\right) =k\quad \Longrightarrow\quad \pi \\[4pt] =4\left(1-\frac13+\frac15-\frac17+-\cdots\right)}
\end{uncoverenv}
\end{minipage}
\EMside{147.9bp}{\uncover<4->{\EMfigure{1.140}{m01-partial-sums}{The partial sums \EMim{S_1}, \EMim{S_2} and \EMim{S_3} of the Fourier series of the square wave (the dashed line is the square wave itself)}}}
\end{frame}

\def\EMtitlesize{24}
\begin{frame}{Examples of Fourier Series (continued)}
\begin{EMbox}{ex}{1}{\textbf{Example 1.5}}
\EMp{Find the Fourier series of the following function:}
\vspace{3bp}
\EMeq{f(x) =x,\qquad -\pi\le x\le\pi,\qquad f(x+2\pi) =f(x)}
\EMgap
\begin{uncoverenv}<2->
\EMeq{a_0 =\frac1{2\pi}\int_{-\pi}^{\pi}f(x)\,dx =0}
\end{uncoverenv}
\EMgap
\begin{uncoverenv}<3->
\EMeq{a_n =\frac1\pi\int_{-\pi}^{\pi}x\cos nx\,dx =\frac1\pi\left[x\left(\frac1n\sin nx\right)-(1)\left(-\frac1{n^2}\cos nx\right)\right]_{-\pi}^{\pi} =\frac1\pi\left[\frac1{n^2}\big(\cos n\pi-\cos(-n\pi)\big)\right] =0}
\end{uncoverenv}
\EMgap
\begin{uncoverenv}<4->
\EMeq{b_n =\frac1\pi\int_{-\pi}^{\pi}x\sin nx\,dx =\frac2\pi\int_0^\pi x\sin nx\,dx =\frac2\pi\left[x\left(\frac{-1}n\cos nx\right)-(1)\left(\frac{-1}{n^2}\sin nx\right)\right]_0^\pi =\frac2\pi\left[\frac{-\pi}n\cos n\pi\right] \\[4pt] =\frac2n(-1)^{n+1}}
\end{uncoverenv}
\EMgap
\begin{uncoverenv}<5->
\EMeq{f(x) =\sum_{n=1}^{\infty}\frac2n(-1)^{n+1}\sin nx =2\left(\sin x-\frac12\sin2x+\frac13\sin3x-+\cdots\right)}
\end{uncoverenv}
\end{EMbox}
\end{frame}

\def\EMtitlesize{24}
\begin{frame}{Examples of Fourier Series (continued)}
\begin{minipage}[t]{500bp}\vspace{0pt}
\EMgap
\begin{uncoverenv}<3->
\EMp{Setting \EMim{x=\frac\pi2}:}
\vspace{3bp}
\EMeq{f\!\left(\frac\pi2\right) =\frac\pi2 =\sum_{n=1}^{\infty}\frac2n(-1)^{n+1}\sin\!\left(n\frac\pi2\right) \\[4pt] =2\left(1-\frac13+\frac15-\frac17+-\cdots\right) =\sum_{n=1}^{\infty}\frac2{2n-1}(-1)^{n+1}\  \Longrightarrow\ \pi \\[4pt] =\sum_{n=1}^{\infty}\frac4{2n-1}(-1)^{n+1}}
\end{uncoverenv}
\EMgap
\begin{uncoverenv}<4->
\EMp{The value of the series at the point of discontinuity \EMim{x=\pi}:}
\vspace{3bp}
\EMeq{f(\pi) =\;?\quad \Longrightarrow\quad \sum_{n=1}^{\infty}\frac2n(-1)^{n+1}\sin(n\pi) =0}
\end{uncoverenv}
\end{minipage}
\EMside{0.0bp}{\uncover<2->{\EMfigure{0.930}{m01-sawtooth}{The sawtooth wave \EMim{f(x)=x} on \EMim{[-\pi,\pi]} with its periodic extension}}}
\end{frame}

\def\EMtitlesize{24}
\begin{frame}{Examples of Fourier Series (continued)}
\begin{EMbox}{imp}{1}{\textbf{Key result 1.6}\ \ |\ \ Important remarks}
\begin{uncoverenv}<2->
\EMitem{1.}{\dimexpr\linewidth-26bp\relax}{If a periodic function is odd, all the \EMim{a} coefficients are zero.}
\end{uncoverenv}
\EMgap
\begin{uncoverenv}<3->
\EMitem{2.}{\dimexpr\linewidth-26bp\relax}{At a point of discontinuity, the value of the Fourier series equals the average of the upper and lower limits at that point.}
\end{uncoverenv}
\end{EMbox}
\end{frame}

\def\EMtitlesize{24}
\begin{frame}{Examples of Fourier Series (continued)}
\begin{EMbox}{ex}{1}{\textbf{Example 1.7}}
\EMp{Find the Fourier series of the following function:}
\vspace{3bp}
\EMeq{f(x) =x^2,\qquad -\pi\le x\le\pi,\qquad f(x+2\pi) =f(x)}
\EMgap
\begin{uncoverenv}<2->
\EMeq{a_0 =\frac1{2\pi}\int_{-\pi}^{\pi}x^2\,dx =\frac{2}{2\pi}\int_0^\pi x^2\,dx =\frac1\pi\left[\frac13x^3\right]_0^\pi =\frac{\pi^2}{3}}
\end{uncoverenv}
\EMgap
\begin{uncoverenv}<3->
\EMeq{a_n =\frac1\pi\int_{-\pi}^{\pi}x^2\cos nx\,dx =\frac2\pi\int_0^\pi x^2\cos nx\,dx =\frac2\pi\left[x^2\left(\frac1n\sin nx\right)-2x\left(\frac{-1}{n^2}\cos nx\right)+2\left(\frac{-1}{n^3}\sin nx\right)\right]_0^\pi \\[4pt] =\frac2\pi\left[\frac{2\pi}{n^2}\cos n\pi\right] =\frac4{n^2}(-1)^n}
\end{uncoverenv}
\EMgap
\begin{uncoverenv}<4->
\EMeq{b_n =\frac1\pi\int_{-\pi}^{\pi}x^2\sin nx\,dx =0}
\end{uncoverenv}
\EMgap
\begin{uncoverenv}<5->
\EMeq{f(x) =\frac{\pi^2}3+\sum_{n=1}^{\infty}\frac4{n^2}(-1)^n\cos nx}
\end{uncoverenv}
\end{EMbox}
\end{frame}

\def\EMtitlesize{24}
\begin{frame}{Examples of Fourier Series (continued)}
\begin{minipage}[t]{500bp}\vspace{0pt}
\EMgap
\begin{EMbox}{imp}{2}{\textbf{Key result 1.8}\ \ |\ \ Important remark}
\begin{uncoverenv}<3->
\EMp{If a periodic function is even, all the \EMim{b} coefficients are zero.}
\end{uncoverenv}
\end{EMbox}
\begin{uncoverenv}<4->
\EMp{Setting \EMim{x=0}:}
\vspace{3bp}
\EMeq{f(0) =0 =\frac{\pi^2}3+\sum_{n=1}^{\infty}\frac4{n^2}(-1)^n\cos(n\cdot0)\quad \Longrightarrow\quad \frac{\pi^2}3 \\[4pt] =\sum_{n=1}^{\infty}\frac4{n^2}(-1)^{n+1}}
\end{uncoverenv}
\EMgap
\begin{uncoverenv}<5->
\EMeq{\pi^2 =\sum_{n=1}^{\infty}\frac{12}{n^2}(-1)^{n+1} =12\left(1-\frac14+\frac19-\frac1{16}+-\cdots\right)}
\end{uncoverenv}
\end{minipage}
\EMside{0.0bp}{\uncover<2->{\EMfigure{1.140}{m01-parabola}{The function \EMim{f(x)=x^2} on \EMim{[-\pi,\pi]} with its periodic extension}}}
\end{frame}

\def\EMtitlesize{24}
\begin{frame}{Examples of Fourier Series (continued)}
\begin{EMbox}{ex}{1}{\textbf{Example 1.9}}
\EMp{Find the Fourier series of the following function:}
\vspace{3bp}
\EMeq{f(x) =x^2,\qquad 0\le x\le2\pi,\qquad f(x+2\pi) =f(x)}
\EMgap
\begin{uncoverenv}<2->
\EMeq{a_0 =\frac1{2\pi}\int_0^{2\pi}x^2\,dx =\frac1{2\pi}\left[\frac13x^3\right]_0^{2\pi} =\frac1{2\pi}\cdot\frac138\pi^3 =\frac{4\pi^2}3}
\end{uncoverenv}
\EMgap
\begin{uncoverenv}<3->
\EMeq{a_n =\frac1\pi\int_0^{2\pi}x^2\cos nx\,dx =\frac1\pi\left[x^2\left(\frac1n\sin nx\right)-2x\left(\frac{-1}{n^2}\cos nx\right)+2\left(\frac{-1}{n^3}\sin nx\right)\right]_0^{2\pi} =\frac1\pi(2\times2\pi)\left(\frac1{n^2}\cos2\pi n\right) \\[4pt] =\frac4{n^2}}
\end{uncoverenv}
\EMgap
\begin{uncoverenv}<4->
\EMeq{b_n =\frac1\pi\int_0^{2\pi}x^2\sin nx\,dx =\frac1\pi\left[x^2\left(\frac{-1}n\cos nx\right)-2x\left(\frac{-1}{n^2}\sin nx\right)+2\left(\frac1{n^3}\cos nx\right)\right]_0^{2\pi} \\[4pt] =\frac1\pi\left[(4\pi^2)\left(\frac{-1}n\right)+2\cdot\frac1{n^3}-2\cdot\frac1{n^3}\right] =-\frac{4\pi}n}
\end{uncoverenv}
\end{EMbox}
\end{frame}

\def\EMtitlesize{24}
\begin{frame}{Examples of Fourier Series (continued)}
\begin{minipage}[t]{500bp}\vspace{0pt}
\begin{EMbox}{ex}{1}{}
\begin{uncoverenv}<2->
\EMeq{f(x) =\frac{4\pi^2}3+\sum_{n=1}^{\infty}\left(\frac4{n^2}\cos nx-\frac{4\pi}n\sin nx\right)}
\end{uncoverenv}
\end{EMbox}
\end{minipage}
\EMside{95.1bp}{\uncover<3->{\EMfigure{1.140}{m01-shifted-parabola}{The function \EMim{f(x)=x^2} on \EMim{[0,2\pi]} with its periodic extension}}}
\end{frame}

\def\EMtitlesize{24}
\begin{frame}{Fourier Series of Functions with an Arbitrary Period}
\begin{uncoverenv}<2->
\EMeq{f(x+T) =f(x)}
\end{uncoverenv}
\EMgap
\begin{uncoverenv}<3->
\EMp{Using the auxiliary variable \EMim{t=\dfrac{2\pi}{T}x} (\EMim{x:0\to T} and \EMim{t:0\to2\pi}):}
\vspace{3bp}
\EMeq{f\!\left(\frac T{2\pi}t\right) =a_0+\sum_{n=1}^{\infty}a_n\cos nt+b_n\sin nt}
\end{uncoverenv}
\EMgap
\begin{uncoverenv}<4->
\EMeq{a_0 =\frac1{2\pi}\int_{2\pi}f\!\left(\frac T{2\pi}t\right)dt,\qquad a_n =\frac1\pi\int_{2\pi}f\!\left(\frac T{2\pi}t\right)\cos nt\,dt,\qquad b_n =\frac1\pi\int_{2\pi}f\!\left(\frac T{2\pi}t\right)\sin nt\,dt}
\end{uncoverenv}
\EMgap
\begin{uncoverenv}<5->
\EMp{Returning to \EMim{x=\dfrac T{2\pi}t}:}
\vspace{3bp}
\EMeq{f(x) =a_0+\sum_{n=1}^{\infty}a_n\cos\frac{2\pi}Tnx+b_n\sin\frac{2\pi}Tnx}
\end{uncoverenv}
\EMgap
\begin{uncoverenv}<6->
\EMeq{a_0 =\frac1{2\pi}\int_Tf(x)\frac{2\pi}T\,dx =\frac1T\int_Tf(x)\,dx}
\end{uncoverenv}
\end{frame}

\def\EMtitlesize{24}
\begin{frame}{Fourier Series of Functions with an Arbitrary Period (continued)}
\begin{uncoverenv}<2->
\EMeq{a_n =\frac1\pi\int_Tf(x)\cos\!\left(\frac{2\pi}Tnx\right)\frac{2\pi}T\,dx =\frac2T\int_Tf(x)\cos\!\left(\frac{2\pi}Tnx\right)dx}
\end{uncoverenv}
\EMgap
\begin{uncoverenv}<3->
\EMeq{b_n =\frac1\pi\int_Tf(x)\sin\!\left(\frac{2\pi}Tnx\right)\frac{2\pi}T\,dx =\frac2T\int_Tf(x)\sin\!\left(\frac{2\pi}Tnx\right)dx}
\end{uncoverenv}
\EMgap
\begin{EMbox}{thm}{3}{\textbf{Theorem 1.10}\ \ |\ \ The Fourier series (period \EMim{T})}
\begin{uncoverenv}<4->
\EMp{If the function \EMim{f(x)} is periodic with period \EMim{T}, it can be expressed as a trigonometric series as follows:}
\vspace{3bp}
\EMeq{f(x) =a_0+\sum_{n=1}^{\infty}a_n\cos\frac{2\pi}Tnx+b_n\sin\frac{2\pi}Tnx,\qquad f(x+T) =f(x)}
\end{uncoverenv}
\EMgap
\begin{uncoverenv}<5->
\EMeq{a_0 =\frac1T\int_Tf(x)\,dx,\qquad a_n =\frac2T\int_Tf(x)\cos\frac{2\pi}Tnx\,dx,\qquad b_n =\frac2T\int_Tf(x)\sin\frac{2\pi}Tnx\,dx}
\end{uncoverenv}
\end{EMbox}
\end{frame}

\def\EMtitlesize{24}
\begin{frame}{Fourier Series of Functions with an Arbitrary Period (continued)}
\begin{EMbox}{rem}{1}{\textbf{Remark}}
\begin{uncoverenv}<2->
\EMp{At points of discontinuity of the function, the value of the Fourier series equals the average of the upper and lower limits there; the value of the Fourier series at the point \EMim{x} is}
\vspace{3bp}
\EMeq{\frac{f(x^+)+f(x^-)}{2}}
\end{uncoverenv}
\end{EMbox}
\end{frame}

\def\EMtitlesize{24}
\begin{frame}{Dirichlet's Theorem}
\begin{minipage}[t]{500bp}\vspace{0pt}
\begin{EMbox}{thm}{1}{\textbf{Theorem 1.11}\ \ |\ \ Dirichlet's theorem}
\begin{uncoverenv}<2->
\EMp{If the function \EMim{f(x)} is periodic and \textbf{bounded}, has a \textbf{finite number of maxima and minima}, and also a \textbf{finite number of points of discontinuity} in each period, then this function has a Fourier series, and the value of the Fourier series at points of discontinuity equals the average of the right and left limits at those points.}
\end{uncoverenv}
\end{EMbox}
\begin{uncoverenv}<3->
\EMp{Examples of functions that do not satisfy Dirichlet's conditions:}
\end{uncoverenv}
\EMgap
\begin{uncoverenv}<4->
\EMeq{f(x) =\frac1x,\quad 0<x\le1,\qquad f(x+1) =f(x)}
\end{uncoverenv}
\EMgap
\begin{uncoverenv}<5->
\EMeq{f(x) =\sin\frac1x,\quad 0<x\le1,\qquad f(x+1) =f(x)}
\end{uncoverenv}
\EMgap
\end{minipage}
\EMside{253.0bp}{\uncover<6->{\EMfigure{0.935}{m01-dirichlet}{Functions that violate Dirichlet's conditions: \EMim{\frac1x} (unbounded), \EMim{\sin\frac1x} (infinitely many maxima and minima) and a staircase function with infinitely many steps}}}
\end{frame}

\def\EMtitlesize{24}
\begin{frame}{Solving Some Problems}
\begin{EMbox}{ex}{1}{\textbf{Example 1.12}\ \ |\ \ Problem 1}
\EMp{Find the Fourier series of the following function:}
\vspace{3bp}
\EMeq{f(x) =\sin\alpha x,\qquad \alpha\in\mathbb{R},\qquad x\in[-\pi,\pi),\qquad f(x+2\pi) =f(x)}
\EMgap
\begin{uncoverenv}<2->
\EMp{Since this function is odd, \EMim{a_0=0} and \EMim{a_n=0}.}
\vspace{3bp}
\EMeq{\begin{aligned}
b_n&=\frac1\pi\int_{-\pi}^{\pi}\sin\alpha x\sin nx\,dx=\frac2\pi\int_0^\pi\sin\alpha x\sin nx\,dx=\frac1\pi\int_0^\pi\big[\cos(\alpha-n)x-\cos(\alpha+n)x\big]dx\\
&=\frac1\pi\left[\frac{\sin(\alpha-n)x}{\alpha-n}-\frac{\sin(\alpha+n)x}{\alpha+n}\right]_0^\pi=\cdots=\frac{2n(-1)^n\sin\alpha\pi}{\pi(\alpha^2-n^2)}
\end{aligned}}
\end{uncoverenv}
\EMgap
\begin{uncoverenv}<3->
\EMeq{f(x) =\sum_{n=1}^{\infty}\frac{2n(-1)^n\sin\alpha\pi}{\pi(\alpha^2-n^2)}\sin nx}
\end{uncoverenv}
\EMgap
\begin{uncoverenv}<4->
\EMp{We examine the series when \EMim{\alpha=1}:}
\vspace{3bp}
\EMeq{\lim_{\alpha\to1}b_1 =\lim_{\alpha\to1}\frac{2(-1)^1\sin\alpha\pi}{\pi(\alpha^2-1^2)} =\lim_{\alpha\to1}\frac{-2(\pi\cos\alpha\pi)}{\pi(2\alpha)} =1,\qquad b_n\Big|_{n\neq1} =0}
\end{uncoverenv}
\end{EMbox}
\end{frame}

\def\EMtitlesize{24}
\begin{frame}{Solving Some Problems (continued)}
\begin{EMbox}{ex}{1}{}
\begin{uncoverenv}<2->
\EMp{Hence the Fourier series is simply \EMim{f(x)=\sin x}.}
\end{uncoverenv}
\end{EMbox}
\end{frame}

\def\EMtitlesize{24}
\begin{frame}{Solving Some Problems (continued)}
\begin{EMbox}{exe}{1}{\textbf{Exercise 1.13}}
\EMp{Examine the series when \EMim{\alpha=m}, \EMim{m\in\mathbb{N}}.}
\end{EMbox}
\end{frame}

\def\EMtitlesize{24}
\begin{frame}{Solving Some Problems (continued)}
\begin{EMbox}{ex}{1}{\textbf{Example 1.14}\ \ |\ \ Problem 2}
\EMp{Find the Fourier series of the following function and deduce \EMim{\csc\alpha\pi} from it.}
\vspace{3bp}
\EMeq{f(x) =\cos\alpha x,\qquad \alpha\in\mathbb{R},\qquad x\in[-\pi,\pi),\qquad f(x+2\pi) =f(x)}
\EMgap
\begin{uncoverenv}<2->
\EMp{Since this function is even, \EMim{b_n=0} and \EMim{f(x)=a_0+\sum_{n=1}^{\infty}a_n\cos nx}.}
\vspace{3bp}
\EMeq{a_0 =\frac1{2\pi}\int_{-\pi}^{\pi}\cos\alpha x\,dx =\frac1\pi\int_0^\pi\cos\alpha x\,dx =\frac1\pi\left[\frac{\sin\alpha x}{\alpha}\right]_0^\pi =\frac1{\alpha\pi}\sin\alpha\pi}
\end{uncoverenv}
\EMgap
\begin{uncoverenv}<3->
\EMeq{a_n =\frac2\pi\int_0^\pi\cos\alpha x\cos nx\,dx =\frac1\pi\int_0^\pi\big[\cos(\alpha-n)x+\cos(\alpha+n)x\big]dx =\frac1\pi\left[\frac{\sin(\alpha-n)x}{\alpha-n}+\frac{\sin(\alpha+n)x}{\alpha+n}\right]_0^\pi =\cdots \\[4pt] =\frac{(-1)^n2\alpha\sin\alpha\pi}{\pi(\alpha^2-n^2)}}
\end{uncoverenv}
\EMgap
\begin{uncoverenv}<4->
\EMeq{f(x) =\frac1{\alpha\pi}\sin\alpha\pi+\frac{2\alpha\sin\alpha\pi}\pi\sum_{n=1}^{\infty}\frac{(-1)^n}{(\alpha^2-n^2)}\cos nx}
\end{uncoverenv}
\end{EMbox}
\end{frame}

\def\EMtitlesize{24}
\begin{frame}{Solving Some Problems (continued)}
\begin{EMbox}{ex}{1}{}
\begin{uncoverenv}<2->
\EMp{The periodic function is continuous at \EMim{x=0}, so}
\vspace{3bp}
\EMeq{f(0) =1 =\frac1{\alpha\pi}\sin\alpha\pi+\frac{2\alpha\sin\alpha\pi}\pi\sum_{n=1}^{\infty}\frac{(-1)^n}{(\alpha^2-n^2)}\quad \Longrightarrow\quad \frac1{\sin\alpha\pi} =\csc\alpha\pi =\frac1{\alpha\pi}+\frac{2\alpha}\pi\sum_{n=1}^{\infty}\frac{(-1)^n}{(\alpha^2-n^2)}}
\end{uncoverenv}
\end{EMbox}
\end{frame}

\def\EMtitlesize{24}
\begin{frame}{Solving Some Problems (continued)}
\begin{EMbox}{exe}{1}{\textbf{Exercise 1.15}}
\EMp{Prove that if \EMim{\alpha=m}, \EMim{m\in\mathbb{N}}, then}
\vspace{3bp}
\EMeq{a_0 =0,\qquad a_n =\begin{cases}1&n=m\\ 0&n\neq m\end{cases}}
\end{EMbox}
\end{frame}

\def\EMtitlesize{24}
\begin{frame}{Solving Some Problems (continued)}
\begin{EMbox}{ex}{1}{\textbf{Example 1.16}\ \ |\ \ Problem 3}
\EMp{Find the Fourier series of the following function and deduce \EMim{\coth\alpha\pi} from it.}
\vspace{3bp}
\EMeq{f(x) =e^{-\alpha x},\qquad \alpha\in\mathbb{R},\qquad x\in[-\pi,\pi),\qquad f(x+2\pi) =f(x)}
\EMgap
\begin{uncoverenv}<2->
\EMp{\textbf{Lemma 1:}}
\vspace{3bp}
\EMeq{\int e^{\alpha x}\cos nx\,dx =\frac{e^{\alpha x}}{\alpha^2+n^2}\big(\alpha\cos nx+n\sin nx\big)}
\end{uncoverenv}
\EMgap
\begin{uncoverenv}<3->
\EMp{Proof: integrating by parts twice, with \EMim{I=\int e^{\alpha x}\cos nx\,dx}:}
\vspace{3bp}
\EMeq{I =\frac1\alpha e^{\alpha x}\cos nx-\int\frac1\alpha e^{\alpha x}(-n\sin nx)\,dx =\frac1\alpha e^{\alpha x}\cos nx+\frac n\alpha\int e^{\alpha x}\sin nx\,dx}
\end{uncoverenv}
\EMgap
\begin{uncoverenv}<4->
\EMeq{=\frac1\alpha e^{\alpha x}\cos nx+\frac n\alpha\left[\frac1\alpha e^{\alpha x}\sin nx-\int\frac1\alpha e^{\alpha x}(n\cos nx)\,dx\right] =\frac1\alpha e^{\alpha x}\cos nx+\frac n{\alpha^2}e^{\alpha x}\sin nx-\frac{n^2}{\alpha^2}I}
\end{uncoverenv}
\EMgap
\begin{uncoverenv}<5->
\EMeq{\Longrightarrow\quad I =\frac{e^{\alpha x}}{\alpha^2+n^2}\big(\alpha\cos nx+n\sin nx\big)}
\end{uncoverenv}
\end{EMbox}
\end{frame}

\def\EMtitlesize{24}
\begin{frame}{Solving Some Problems (continued)}
\begin{EMbox}{ex}{1}{}
\begin{uncoverenv}<2->
\EMp{\textbf{Lemma 2:}}
\vspace{3bp}
\EMeq{\int e^{\alpha x}\sin nx\,dx =\frac{e^{\alpha x}}{\alpha^2+n^2}\big(\alpha\sin nx-n\cos nx\big)}
\end{uncoverenv}
\EMgap
\begin{uncoverenv}<3->
\EMp{(proof left as an exercise.)}
\end{uncoverenv}
\EMgap
\begin{uncoverenv}<4->
\EMp{The function is neither even nor odd, so all the Fourier coefficients must be computed.}
\vspace{3bp}
\EMeq{a_0 =\frac1{2\pi}\int_{-\pi}^{\pi}e^{-\alpha x}\,dx =\frac1{-2\pi\alpha}\Big[e^{-\alpha x}\Big]_{-\pi}^{\pi} =\frac1{2\pi\alpha}\big[e^{\alpha\pi}-e^{-\alpha\pi}\big] =\frac{\sinh\alpha\pi}{\alpha\pi}}
\end{uncoverenv}
\EMgap
\begin{uncoverenv}<5->
\EMeq{a_n =\frac1\pi\int_{-\pi}^{\pi}e^{-\alpha x}\cos nx\,dx =\left\{\frac{e^{-\alpha x}}{\pi(\alpha^2+n^2)}\big(-\alpha\cos nx+n\sin nx\big)\right\}_{-\pi}^{\pi} =\frac{\alpha(-1)^n}{\pi(\alpha^2+n^2)}\big[e^{\alpha\pi}-e^{-\alpha\pi}\big] \\[4pt] =\frac{2\alpha(-1)^n}{\pi(\alpha^2+n^2)}\sinh\alpha\pi}
\end{uncoverenv}
\end{EMbox}
\end{frame}

\def\EMtitlesize{24}
\begin{frame}{Solving Some Problems (continued)}
\begin{EMbox}{ex}{1}{}
\begin{uncoverenv}<2->
\EMeq{b_n =\frac1\pi\int_{-\pi}^{\pi}e^{-\alpha x}\sin nx\,dx =\left\{\frac{e^{-\alpha x}}{\pi(\alpha^2+n^2)}\big(-\alpha\sin nx-n\cos nx\big)\right\}_{-\pi}^{\pi} =\frac{n(-1)^n}{\pi(\alpha^2+n^2)}\big[e^{\alpha\pi}-e^{-\alpha\pi}\big] \\[4pt] =\frac{2n(-1)^n}{\pi(\alpha^2+n^2)}\sinh\alpha\pi}
\end{uncoverenv}
\EMgap
\begin{uncoverenv}<3->
\EMeq{\begin{aligned}
f(x)&=\frac{\sinh\alpha\pi}{\alpha\pi}+\sum_{n=1}^{\infty}\left(\frac{2\alpha(-1)^n}{\pi(\alpha^2+n^2)}\sinh\alpha\pi\cos nx+\frac{2n(-1)^n}{\pi(\alpha^2+n^2)}\sinh\alpha\pi\sin nx\right)\\
&=\sinh\alpha\pi\left(\frac1{\alpha\pi}+\sum_{n=1}^{\infty}\left(\frac{2\alpha(-1)^n}{\pi(\alpha^2+n^2)}\cos nx+\frac{2n(-1)^n}{\pi(\alpha^2+n^2)}\sin nx\right)\right)
\end{aligned}}
\end{uncoverenv}
\EMgap
\begin{uncoverenv}<4->
\EMp{Since \EMim{x=\pi} is a point of discontinuity, the value of the Fourier series there equals the average of the upper and lower limits of the function at that point, i.e.}
\vspace{3bp}
\EMeq{\frac{f(\pi^+)+f(\pi^-)}2 =\frac{e^{\alpha\pi}+e^{-\alpha\pi}}2 =\cosh\alpha\pi}
\end{uncoverenv}
\end{EMbox}
\end{frame}

\def\EMtitlesize{24}
\begin{frame}{Solving Some Problems (continued)}
\begin{EMbox}{ex}{1}{}
\begin{uncoverenv}<2->
\EMeq{\cosh\alpha\pi =\sinh\alpha\pi\left(\frac1{\alpha\pi}+\sum_{n=1}^{\infty}\left(\frac{2\alpha(-1)^n}{\pi(\alpha^2+n^2)}\cos n\pi+\frac{2n(-1)^n}{\pi(\alpha^2+n^2)}\sin n\pi\right)\right) =\sinh\alpha\pi\left(\frac1{\alpha\pi}+\sum_{n=1}^{\infty}\frac{2\alpha}{\pi(\alpha^2+n^2)}\right)}
\end{uncoverenv}
\EMgap
\begin{uncoverenv}<3->
\EMeq{\coth\alpha\pi =\frac{\cosh\alpha\pi}{\sinh\alpha\pi} =\frac1{\alpha\pi}+\sum_{n=1}^{\infty}\frac{2\alpha}{\pi(\alpha^2+n^2)}}
\end{uncoverenv}
\end{EMbox}
\end{frame}
